% !TeX root = ../se200_referential_regimes.tex

% =============================================================================
\section{Nine Candidate Carrier--Basis Profiles}
\label{sec:nine}
% =============================================================================

The selected construction of Sections~\ref{sec:inventory} and
\ref{sec:algebra} contains nine named carrier--basis profiles.
Three carrier kinds are represented by one candidate basis each,
and three by pairs of candidate bases.
This is a catalog selected for analysis, not a claim that all nine
profiles are scientifically necessary or irreducible.

The following labels name the selected profiles:

\begin{description}[style=nextline,leftmargin=2.5cm]
  \item[$\OBL$] obligation-bearer identity.
  \item[$\OCC$] occurrence identity.
  \item[$\REC$] record identity.
  \item[$\LOC$] locus-fixed plain-referent identity.
  \item[$\OBJ$] object-fixed plain-referent identity.
  \item[$\SCOPEE$] extension-fixed scope identity.
  \item[$\SCOPES$] structure-fixed scope identity.
  \item[$\RULEC$] content-fixed rule identity.
  \item[$\RULES$] structure-fixed rule identity.
\end{description}

Table~\ref{tab:master} records proposed classifications over the finite
transformation basis $\Fset$.
Its entries are stipulated family-level profile assignments for the
selected comparison settings.
They must not be read as established
identity effects of every operation in the named family or as
independently warranted judgments about real-world referents.
Where eligible operations within a family have different effects,
that family requires a more specific classification domain.

\begin{table}[ht]
  \centering
  \small
  \setlength{\tabcolsep}{4pt}
  \begin{tabular}{@{}l ccc cc cc cc@{}}
    \toprule
       & OBL  & OCC  & REC  & LOC  & OBJ  & SCOPE-E & SCOPE-S & RULE-C & RULE-S \\
    \midrule
    RE & \NEU & \NEU & \NEU & \NEU & \NEU & \NEU    & \NEU    & \NEU   & \NEU   \\
    AN & \NEU & \NEU & \PRS & \NEU & \NEU & \NEU    & \NEU    & \NEU   & \NEU   \\
    RF & \PRS & \PRS & \PRS & \PRS & \PRS & \PRS    & \BRK    & \PRS   & \BRK   \\
    AD & \PRS & \BRK & \PRS & \PRS & \PRS & \PRS    & \BRK    & \PRS   & \BRK   \\
    RC & \NEU & \NEU & \NEU & \NEU & \NEU & \PRS    & \BRK    & \NEU   & \NEU   \\
    RA & \NEU & \NA  & \NEU & \NA  & \NA  & \NEU    & \BRK    & \NEU   & \NEU   \\
    SU & \BRK & \BRK & \BRK & \BRK & \BRK & \BRK    & \BRK    & \BRK   & \BRK   \\
    BF & \NEU & \NA  & \PRS & \PRS & \BRK & \PRS    & \BRK    & \BRK   & \BRK   \\
    PV & \NEU & \NEU & \PRS & \NEU & \NEU & \NEU    & \NEU    & \NEU   & \NEU   \\
    SE & \NEU & \PRS & \NEU & \PRS & \PRS & \NA     & \NA     & \NA    & \NA    \\
    \bottomrule
  \end{tabular}
  \caption{Proposed classification profiles over $\Fset$ for the
    selected carrier--basis construction.
    \PRS{}: preservation of the declared basis;
    \BRK{}: breakage of that basis;
    \NEU{}: applicable but identity-neutral under that basis;
    \NA{}: inapplicable on the selected domain.
    Entries are candidate assignments, not universal identity laws.}
  \label{tab:master}
\end{table}

A \PRS{} entry proposes preservation of the declared identity basis;
a \BRK{} entry proposes its breakage.
An \NEU{} entry indicates an applicable operation that does not affect
the basis, and an \NA{} entry indicates inapplicability on the selected
domain.
Neither \PRS{} nor \NEU{} independently warrants strict identity
between actual referents.
In particular, a \BRK{} entry in a proposed profile does not by itself
establish that two specified objects are distinct.

For the selected same-carrier comparisons, BF distinguishes the
proposed $\LOC$ and $\OBJ$ classifications; AD distinguishes
$\SCOPEE$ and $\SCOPES$; and RF distinguishes $\RULEC$
and $\RULES$.
The conditional examples in Section~\ref{sec:algebra} show how
stipulated locus, object, extension, structure, and content criteria
can yield these classification differences.
The other entries record additional proposed behavior whose adequacy
also requires examination; they are not needed to establish the
three conditional split comparisons.

The uniform RE and SU rows are features of this selected table:
RE is assigned \NEU{} throughout, while SU is assigned \BRK{} throughout.
Those assignments are not general theorems about re-expression or
substitution under every scientifically adequate identity criterion.

A disagreement between two classification functions is distinct from
a disagreement between the equivalence relations generated by
non-breaking steps.
To state the latter, the two relations must be
compared on a common eligible domain.

\begin{definition}[Separation witness]
  \label{se200.def.SeparationWitness}
  Let $\tau$ and $\tau'$ be profiles with a common domain $D$ of
  eligible presentations.
  A \emph{separation witness for their generated
  relations on $D$} is a pair $(r,s)\in D\times D$ for which
  \[
    r\sim_\tau s \quad\text{and}\quad r\not\sim_{\tau'}s,
  \]
  or conversely, with each relation restricted to $D$.
\end{definition}

Such a witness demonstrates that the generated relations differ on
$D$; it does not independently establish the correctness of either
relation as strict real-world identity.
Nor does a difference in the
classification of a transformation family necessarily provide such
a witness, since equivalence closure may connect the relevant pair
by another non-breaking path.

\begin{proposition}[Distinct core profiles]
  \label{se200.prop.DistinctCoreProfiles}
  The nine classification functions specified by the columns of
  Table~\ref{tab:master} are pairwise distinct.
\end{proposition}

\begin{proof}
  It is enough to inspect the entries at AN, RF, RA, and BF.
  The four-entry classification signatures in these rows are
  pairwise different across all nine columns of
  Table~\ref{tab:master}.
  Hence no two of the nine
  classification functions agree on all of $\Fset$.
\end{proof}

This proposition establishes distinctness of nine \emph{declared
classification functions} in the selected catalog.
It does not establish pairwise distinctness of their generated
relations on a common domain, that all nine identity bases are
independently warranted, or that any faithful realization requires
nine non-amalgamable regimes.
Those questions are examined as limits of the counting argument in
Section~\ref{sec:bound}.
